\subsection{Automatic Software Documentation Generation}
Automatic software documentation generation methods can be categorized into three types. Template-based methods~\cite{template2, template3, template4} parse specific information from source code and then populate it into predefined templates. For instance, Hill et al.~\cite{template1} generate annotations by analyzing the identifiers of Java methods. Information retrieval-based methods retrieve suitable descriptions from a vast documentation corpus~\cite{info-retrieve1, info-retrieve4, info-retrieve5}, including bug tracking systems~\cite{info-retrieve2} and developer forums like Stack Overflow~\cite{info-retrieve3}. Deep learning-based methods represent a major focus of current research~\cite{deep-learning1, deep-learning2, deep-learning5}. DocAgent~\cite{docagent} designs a multi-agent framework to generate high-quality documentation. RepoAgent~\cite{repoagent} utilizes global context to infer code functionality and semantics. 

\subsection{Software Documentation Evaluation}
Existing evaluation methods for software documentation can be classified into three categories. Human-based methods~\cite{human1, human2, human3} invite experts to provide detailed assessment, which is labor-intensive. Metrics-based methods~\cite{bleu, meteor, cider, rouge} borrow metrics from Natural Language Processing (NLP), focusing on quantifying the textual similarity between the generated documentation and the references. However, these methods typically rely on high-quality reference documentation, which is quite challenging to construct. Nowadays, LLM-as-a-judge methods have gained traction~\cite{hierarchical, docagent} by leveraging the contextual understanding and instruction-following capabilities of LLMs. By providing LLMs with evaluation criteria, they can conduct assessments across different dimensions. 